\chapter{Eliminacja kwantyfikatorów}
\label{ch:mod-qe}
\section{Motywacja i definicja}
Formuła $\exists y\,(x<y\land y<z)$ wyraża istnienie
elementu pomiędzy $x$ i $z$. W gęstym porządku liniowym
jest równoważna formule $x<z$: kwantyfikator nie niesie
tam dodatkowej informacji. Eliminacja kwantyfikatorów
oznacza, że w ustalonej teorii każdą formułę można w ten
sposób zastąpić formułą bez kwantyfikatorów. Nie oznacza
to, że każda formuła jest prawdziwa w każdym modelu.

Pracujemy w logice pierwszego rzędu z równością i w ZFC.
Korzystamy z pojęć termu, formuły i konsekwencji semantycznej
z definicji \ref{def:syntax-term}, \ref{def:syntax-formula}
i \ref{def:sem-entailment}.
Dla wygody dopuszczamy logiczne stałe $\top,\bot$,
o wartościach odpowiednio prawda i fałsz w każdej strukturze.
Jest to konserwatywne rozszerzenie składni: w obecności
zmiennej $x$ można je zastąpić przez $x=x$ i $\neg(x=x)$,
a w obecności termu domkniętego $t$ przez $t=t$ i $\neg(t=t)$.
W języku bez stałych i bez symboli funkcyjnych o arności zero
konwencja ta pozwala mówić także o zdaniach bez kwantyfikatorów.
Jeśli nie ma dostępnych termów domkniętych, nie utożsamiamy
$\top$ i $\bot$ z atomami języka pierwotnego.

\begin{definicja}[Eliminacja kwantyfikatorów]
\label{def:mod-qe}
Niech $T$ będzie spełnialną teorią języka $\mathcal L$.
Mówimy, że $T$ ma eliminację kwantyfikatorów, jeżeli dla
każdej formuły $\varphi(\bar x)$ istnieje formuła
bez kwantyfikatorów $\psi(\bar x)$, bez dodatkowych
zmiennych wolnych, taka że
\[
T\models\forall\bar x\,(\varphi(\bar x)\leftrightarrow\psi(\bar x)).
\]
Dopuszczamy pustą listę $\bar x$, a wówczas uniwersalne
domknięcie oznacza samo zdanie.
\end{definicja}

Równoważność dotyczy wszystkich modeli $T$ i wszystkich
wartościowań wskazanych zmiennych; nie wystarcza równoważność
w jednym wybranym modelu niezupełnej teorii.
Formuła bez kwantyfikatorów jest kombinacją boolowską
formuł atomowych i stałych logicznych.

\begin{lemat}[Eliminacja pojedynczego kwantyfikatora]
\label{lem:mod-qe-exists}
W celu wykazania eliminacji kwantyfikatorów wystarczy
wykazać, że każda formuła postaci $\exists y\,\theta(y,\bar x)$,
gdzie $\theta$ nie zawiera kwantyfikatorów, jest równoważna
modulo $T$ formule bez kwantyfikatorów.
Wystarczy rozpatrywać $\theta$ będące skończonymi
koniunkcjami literałów, tj. atomów i negacji atomów.
\end{lemat}
\begin{proof}
Indukcja po budowie formuły (twierdzenie
\ref{thm:syntax-structural-induction-xlv}) pozostawia atomy
bez zmian i zastępuje podformuły przy spójnikach ich
równoważnikami bez kwantyfikatorów. W kroku egzystencjalnym
stosujemy założenie do już uproszczonej podformuły.
Krok uniwersalny sprowadzamy do negacji i kwantyfikatora
egzystencjalnego. Równoważność podformuł jest zachowywana
przez spójniki i kwantyfikatory z ich semantyki.
Każdą skończoną kombinację boolowską atomów można
zapisać jako alternatywę koniunkcji literałów: dla każdego
wartościowania atomów czyniącego kombinację prawdziwą
tworzymy koniunkcję tych atomów lub ich negacji zgodnie
z wartościowaniem. Alternatywa tych koniunkcji ma tę samą
tablicę prawdy. Pusta koniunkcja to $\top$, pusta
alternatywa to $\bot$. Kwantyfikator egzystencjalny
rozkłada się na skończoną alternatywę, ponieważ świadek
spełnia alternatywę dokładnie wtedy, gdy spełnia
co najmniej jeden jej człon.
\end{proof}

\section{Kryterium semantyczne}
\begin{definicja}[Typ bez kwantyfikatorów]
\label{def:mod-qe-qftype}
Dla krotki $\bar a$ w strukturze $\mathfrak M$ określamy
jej typ bez kwantyfikatorów
$\operatorname{qftp}^{\mathfrak M}(\bar a)$ jako zbiór
wszystkich formuł bez kwantyfikatorów, prawdziwych przy
$\bar a$ i o zmiennych wolnych w ustalonej liście $\bar x$.
Dla pustej krotki jest to zbiór prawdziwych zdań
bez kwantyfikatorów. Typ zupełny $\operatorname{tp}$
definiowano w rozdziale \ref{ch:mod-types}.
\end{definicja}

\begin{twierdzenie}[Kryterium przez typy bez kwantyfikatorów]
\label{thm:mod-qe-semantic}
Dla spełnialnej teorii $T$ równoważne są warunki:
\begin{enumerate}[label=(\roman*)]
\item $T$ ma eliminację kwantyfikatorów;
\item dla dowolnych modeli $\mathfrak M,\mathfrak N\models T$
i skończonych krotek $\bar a,\bar b$ tej samej długości
(w tym pustych) równość ich typów bez kwantyfikatorów
implikuje zgodność prawdziwości wszystkich formuł
przy tych krotkach.
\end{enumerate}
\end{twierdzenie}
\begin{proof}
Jeżeli zachodzi (i), to dowolna formuła ma równoważnik
bez kwantyfikatorów, co daje (ii).

Załóżmy (ii) i ustalmy formułę $\varphi(\bar x)$.
Dla każdej jej realizacji $\bar a$ w pewnym modelu $T$
typ bez kwantyfikatorów $q$ tej krotki wymusza
$\varphi$: inaczej istniałaby realizacja $q$ w modelu $T$,
która nie spełnia $\varphi$, a ponieważ $q$ rozstrzyga
każdą formułę bez kwantyfikatorów, przeczyłoby to (ii).
Formalnie, po dodaniu świeżych stałych $\bar c$,
teoria
\[
T\cup\{\eta(\bar c):\eta\in q\}
  \cup\{\neg\varphi(\bar c)\}
\]
nie ma modelu. Z twierdzenia o zwartości
\ref{thm:compactness} istnieje skończony fragment
$q_0\subseteq q$, którego koniunkcja $\theta_q$ wymusza
$\varphi$ modulo $T$. Podstawienie stałych za zmienne
jest zgodne z semantyką na mocy twierdzenia
\ref{thm:sem-substitution}; dla pustej krotki stałych
nie dodajemy. Koniunkcja $\theta_q$ jest bez kwantyfikatorów
i spełniona przy wyjściowej krotce.

Nie musimy wybierać realizacji ze zbioru wszystkich modeli:
niech $\Theta$ będzie zbiorem wszystkich formuł bez
kwantyfikatorów o wskazanych zmiennych wolnych, które
implikują $\varphi$ modulo $T$. Z poprzedniego akapitu
każda realizacja $\varphi$ spełnia pewną formułę z $\Theta$.
Zatem teoria
\[
T\cup\{\varphi(\bar c)\}
  \cup\{\neg\theta(\bar c):\theta\in\Theta\}
\]
jest niespełnialna. Zwartość daje skończenie wiele formuł
$\theta_1,\ldots,\theta_m$ z $\Theta$, takich że
$\varphi$ implikuje ich alternatywę modulo $T$.
Każdy człon alternatywy implikuje $\varphi$, więc
$\varphi$ jest równoważna tej alternatywie. Jeśli
$m=0$, otrzymujemy $\bot$. Dowód obejmuje także zdania.
\end{proof}

\begin{twierdzenie}[Kryterium przedłużania zgodnych krotek]
\label{thm:mod-qe-forth}
Załóżmy, że dla każdych modeli $\mathfrak M,\mathfrak N$
teorii $T$ i krotek $\bar a,\bar b$ o tym samym typie
bez kwantyfikatorów spełniony jest warunek:
dla każdego $a\in M$ istnieje rozszerzenie elementarne
$\mathfrak N'\succeq\mathfrak N$ i $b\in N'$ takie, że
\[
\operatorname{qftp}^{\mathfrak M}(\bar a,a)=
\operatorname{qftp}^{\mathfrak N'}(\bar b,b).
\]
Wówczas $T$ ma eliminację kwantyfikatorów.
\end{twierdzenie}
\begin{proof}
Weryfikujemy kryterium \ref{thm:mod-qe-semantic}
indukcją po formułach, jednocześnie dla wszystkich
modeli $T$ i zgodnych krotek. Atomy są zgodne z założenia,
a spójniki zachowują zgodność przez indukcję.
Jeśli $\mathfrak M\models\exists y\,\varphi(\bar a,y)$,
wybieramy świadka $a$ i stosujemy założenie.
Hipoteza indukcyjna daje
$\mathfrak N'\models\varphi(\bar b,b)$, stąd
$\mathfrak N'\models\exists y\,\varphi(\bar b,y)$.
Z elementarności rozszerzenia
(\ref{def:model-elementary-substructure})
wynika $\mathfrak N\models\exists y\,\varphi(\bar b,y)$.
Implikacja odwrotna wynika z zastosowania tego samego
założenia po zamianie modeli i krotek. Kwantyfikator
uniwersalny sprowadzamy do negacji i egzystencjalnego.
\end{proof}

\begin{wniosek}[Kryterium w modelach nasyconych]
\label{cor:mod-qe-saturated}
W twierdzeniu \ref{thm:mod-qe-forth} wystarczy sprawdzić
możliwość przedłużenia krotek bez wychodzenia poza modele,
gdy oba modele są $\aleph_0$-nasycone.
\end{wniosek}
\begin{proof}
Dla dowolnych $\mathfrak M,\mathfrak N\models T$ bierzemy
ich $\aleph_0$-nasycone rozszerzenia elementarne
$\mathfrak M',\mathfrak N'$ z twierdzenia
\ref{thm:mod-saturation-extension}. Elementarność
zachowuje typy bez kwantyfikatorów. Stosujemy założenie
do $\bar a,a$ w $\mathfrak M'$ i $\bar b$ w $\mathfrak N'$.
Otrzymujemy świadka w $\mathfrak N'$; jego typ bez
kwantyfikatorów z $\bar b$ zgadza się z typem
$\bar a,a$ także w $\mathfrak M$. Spełnione jest
kryterium \ref{thm:mod-qe-forth}.
\end{proof}

\section{Konsekwencje eliminacji}
\begin{wniosek}[Zanurzenia i zupełność]
\label{cor:mod-qe-consequences}
Jeżeli $T$ ma eliminację kwantyfikatorów, to każde
zanurzenie struktur będących modelami $T$ jest elementarne.
Jeżeli dodatkowo wszystkie modele $T$ zgadzają się co do
zdań bez kwantyfikatorów, to $T$ jest zupełna.
\end{wniosek}
\begin{proof}
Zanurzenie utożsamia dziedzinę z podstrukturą obrazu.
Wartości termów w podstrukturze są takie same jak
w strukturze większej, co wynika z indukcji po termach:
dla zmiennej i stałej jest to definicja podstruktury,
a dla zastosowania symbolu funkcyjnego korzystamy z
ograniczenia interpretacji tego symbolu. Relacje na
podstrukturze są ograniczeniami relacji większej struktury.
Stąd atomy, a przez spójniki wszystkie formuły bez
kwantyfikatorów, mają tę samą wartość logiczną.
Eliminacja kwantyfikatorów rozszerza zgodność na wszystkie
formuły, czyli zanurzenie jest elementarne.
W drugiej części każde zdanie ma równoważnik bez
kwantyfikatorów, a więc wszystkie modele $T$ zgadzają
się co do jego prawdziwości. To jest zupełność teorii.
\end{proof}

\begin{wniosek}[Zbiory definiowalne i typy z parametrami]
\label{cor:mod-qe-definability}
Niech $T$ ma eliminację kwantyfikatorów i
$\mathfrak M\models T$. Każdy zbiór definiowalny
w $M^n$ z parametrami z $A\subseteq M$ jest definiowany
przez formułę bez kwantyfikatorów z parametrami z $A$.
Typ zupełny nad $A$ jest wyznaczony przez należące
do niego formuły bez kwantyfikatorów.
\end{wniosek}
\begin{proof}
Formułę z nazwami skończonej krotki parametrów zapisujemy
jako $\varphi(\bar x,\bar a)$, zgodnie z konwencją
z rozdziału \ref{ch:mod-types}. Stosujemy eliminację
do formuły języka $\mathcal L$ z dodatkową listą
wolnych zmiennych dla parametrów i podstawiamy
ponownie ich nazwy. Równoważność zachodzi przy każdym
wartościowaniu, w szczególności przy wskazanych
parametrach. Daje to pierwszą część.
Dla typu zupełnego nad $A$ wybieramy jego realizację
w rozszerzeniu elementarnym (twierdzenie
\ref{thm:mod-type-elementary-extension}).
Równoważność modulo $T$ przy parametrach z $A$
obowiązuje również w tym rozszerzeniu, więc przynależność
formuły do typu jest równoważna przynależności jej
równoważnika bez kwantyfikatorów.
\end{proof}

\section{Język samej równości i gęste porządki}
\begin{twierdzenie}[Teoria nieskończonych zbiorów]
\label{thm:mod-qe-equality}
Teoria nieskończonych struktur języka samej równości
ma eliminację kwantyfikatorów i jest zupełna.
\end{twierdzenie}
\begin{proof}
Niech $\bar a,\bar b$ będą krotkami w nieskończonych
strukturach o tym samym typie bez kwantyfikatorów.
Atomy są wyłącznie równościami zmiennych, więc
zgadzają się wszystkie równości między współrzędnymi.
Jeżeli $a$ jest jedną ze współrzędnych $\bar a$,
wybieramy odpowiadającą współrzędną $\bar b$ jako $b$.
Zgodność równości sprawia, że wybór nie zależy od
indeksu tej współrzędnej. Jeśli $a$ jest różny od
wszystkich współrzędnych, wybieramy $b$ poza
skończonym zbiorem współrzędnych $\bar b$, co jest
możliwe przez nieskończoność. W obu przypadkach
rozszerzone krotki zgadzają się co do atomów, a więc
i wszystkich formuł bez kwantyfikatorów.
Kryterium \ref{thm:mod-qe-forth} daje eliminację.
W języku nie ma termów domkniętych, a zdania bez
kwantyfikatorów są kombinacjami $\top,\bot$.
Wniosek \ref{cor:mod-qe-consequences} daje zupełność.
\end{proof}

\begin{twierdzenie}[Gęste porządki liniowe bez elementów krańcowych]
\label{thm:mod-qe-dlo}
Teoria gęstych porządków liniowych bez elementu
najmniejszego i największego w języku $\{<\}$
ma eliminację kwantyfikatorów i jest zupełna.
\end{twierdzenie}
\begin{proof}
Przez gęstość rozumiemy warunek
$\forall u\,\forall v\,(u<v\to\exists w\,(u<w\land w<v))$.
Niech $\mathfrak M,\mathfrak N$ będą modelami teorii,
a $\bar a,\bar b$ krotkami zgodnymi co do typu bez
kwantyfikatorów. Równości i porządek między
współrzędnymi są takie same. Porządkujemy różne
wartości współrzędnych $\bar a$ w rosnący ciąg
$u_1<\cdots<u_k$; odpowiadają im różne wartości
$v_1<\cdots<v_k$ w $\bar b$.
Dla nowego elementu $a\in M$ rozważamy przypadki:
\begin{enumerate}
\item Jeżeli $a=u_i$, bierzemy $b=v_i$.
\item Jeżeli $u_i<a<u_{i+1}$, z gęstości wybieramy
$b$ spełniający $v_i<b<v_{i+1}$.
\item Jeżeli $a<u_1$ lub $a>u_k$, wybieramy element
poniżej $v_1$ lub powyżej $v_k$ przez brak elementów
krańcowych.
\item Jeżeli $k=0$, wybieramy dowolny element $N$,
gdyż uniwersum jest niepuste.
\end{enumerate}
Przypadki wyczerpują możliwości przez liniowość porządku.
W każdym z nich wszystkie równości i nierówności
między rozszerzonymi krotkami są zgodne; w tym języku
są to wszystkie atomy. Kryterium
\ref{thm:mod-qe-forth} daje eliminację kwantyfikatorów.
Jak w dowodzie twierdzenia \ref{thm:mod-qe-equality},
wszystkie zdania bez kwantyfikatorów mają ustaloną
wartość, więc teoria jest zupełna.
\end{proof}

\begin{figure}[htbp]
\centering
\begin{tikzpicture}[x=1cm,y=1cm,>=Stealth]
\draw[<->] (0,1.9)--(10,1.9);
\draw[<->] (0,0)--(10,0);
\foreach \x/\i in {1.5/1,4/2,8/3}{
 \fill (\x,1.9) circle (2pt) node[above] {$u_{\i}$};
 \fill (\x,0) circle (2pt) node[below] {$v_{\i}$};
 \draw[densely dashed,gray] (\x,1.65)--(\x,0.25);
}
\fill[blue!65!black] (5.2,1.9) circle (3pt) node[above] {$a$};
\fill[blue!65!black] (6.6,0) circle (3pt) node[below] {$b$};
\draw[->,blue!65!black,thick] (5.3,1.65)--(6.5,0.25);
\node[left] at (0,1.9) {$M$};
\node[left] at (0,0) {$N$};
\end{tikzpicture}
\caption{Przedłużenie zgodnych krotek w gęstych porządkach.
Ważne jest położenie względem wybranych elementów, nie odległość.}
\label{fig:mod-qe-dlo}
\end{figure}

Rysunek \ref{fig:mod-qe-dlo} ilustruje drugi przypadek
dowodu twierdzenia \ref{thm:mod-qe-dlo}.

\begin{przyklad}[Zbiory definiowalne w jednej zmiennej]
\label{ex:mod-qe-dlo-definable}
W gęstym porządku liniowym bez elementów krańcowych
każdy zbiór definiowalny w jednej zmiennej z parametrami
jest skończoną sumą punktów i otwartych przedziałów
(o końcach w zbiorze parametrów lub nieskończonych).
Istotnie, po eliminacji kwantyfikatorów formuła używa
tylko skończenie wielu parametrów. Po uporządkowaniu
ich różnych wartości $a_1<\cdots<a_k$ dzielimy
uniwersum na punkty $a_i$, przedziały $(a_i,a_{i+1})$
i dwa przedziały nieograniczone. Gdy $k=0$, podział
ma jeden człon: całe uniwersum.
Każdy atom z jednym argumentem zmiennym ma na każdym
członie stałą wartość logiczną, więc tak samo ma ją
formuła bez kwantyfikatorów. Zbiór definiowany jest
sumą tych członów, na których formuła jest prawdziwa.
\end{przyklad}

\section{Ciała algebraicznie domknięte}
Pracujemy w języku pierścieni $\{0,1,+,-,\cdot\}$
(z jednoargumentowym minusem). Każdy term tego języka
jest reprezentowany przez wielomian o współczynnikach
całkowitych w swoich zmiennych: dla zmiennych i stałych
jest to oczywiste, a operacje na termach odpowiadają
dodawaniu, negowaniu i mnożeniu wielomianów.
Równość termów można więc zapisać jako $P=0$,
przenosząc prawą stronę na lewą.

\begin{lemat}[Skończone rozgałęzienia algorytmu wielomianowego]
\label{lem:mod-qe-polynomial-branches}
Dla skończonej listy wielomianów w $K[Y]$, których
współczynniki są wielomianami w krotce parametrów
$\bar x$, następujące warunki można wyrazić formułami
bez kwantyfikatorów w języku pierścieni, jednolicie
dla wszystkich ciał $K$:
\begin{enumerate}
\item dany wielomian jest zerowy lub ma wskazany stopień;
\item największy wspólny dzielnik niezerowych
wielomianów z listy ma dodatni stopień;
\item gdy na liście jest niezerowy wielomian,
największy wspólny dzielnik listy dzieli wskazany
wielomian o współczynnikach wielomianowych w $\bar x$.
\end{enumerate}
\end{lemat}
\begin{proof}
Warunek zerowości to koniunkcja równości wszystkich
współczynników z zerem, a warunek stopnia $d$ wymaga
niezerowości współczynnika przy $Y^d$ i zerowości
wszystkich współczynników wyższych. Stopnie są
ograniczone przez stopnie wyjściowych wielomianów.

Na każdej gałęzi rozróżniamy rzeczywiste stopnie
wielomianów, odrzucamy wielomiany zerowe i wykonujemy
algorytm Euklidesa. Dzielenie z resztą polega na
odejmowaniu wielokrotności dzielnika tak, aby
skasować aktualny wyraz wiodący. Współczynniki
wyników są ilorazami wielomianów w $\bar x$:
dzielenie zachodzi jedynie przez współczynnik wiodący,
którego niezerowość jest już warunkiem danej gałęzi.
Rozróżniamy stopień każdej reszty, sprawdzając
zerowość lub niezerowość liczników tych ilorazów
(mianowniki są niezerowe na tej gałęzi).
Każde dzielenie ma skończenie wiele kroków, ograniczonych
przez stopień dzielnej plus jeden, a kolejne niezerowe
reszty w algorytmie mają ściśle malejące stopnie.
Wobec ustalonych początkowych ograniczeń stopni i
skończonej listy wielomianów drzewo możliwych rozgałęzień
jest skończone. Największy wspólny dzielnik otrzymujemy
iterując algorytm dla listy: dla pary $f,g$ z $g\ne0$
wspólne dzielniki $f,g$ to dokładnie wspólne dzielniki
$g,r$, gdzie $f=qg+r$; ostatnia niezerowa reszta dzieli
wszystkie poprzednie i jest podzielna przez każdy
wspólny dzielnik. Jednoargumentowa lista ma jako
dzielnik swój jedyny niezerowy wielomian.

Na liściu drzewa znamy stopień dzielnika i jego
współczynniki jako ilorazy wielomianów. W przypadku (iii)
dzielimy wskazany wielomian przez ten dzielnik
i sprawdzamy zerowość reszty, ponownie przez
zerowość liczników. Każdy liść jest opisany skończoną
koniunkcją warunków wielomianowych; alternatywa
liści o żądanej własności daje formułę bez kwantyfikatorów.
Nie korzystamy z żadnego założenia o charakterystyce
ciała, ponieważ wszystkie potrzebne niezerowości
sprawdzamy jawnie.
\end{proof}

\begin{twierdzenie}[Eliminacja kwantyfikatorów dla ciał algebraicznie domkniętych]
\label{thm:mod-qe-acf}
Teoria ciał algebraicznie domkniętych w języku
pierścieni ma eliminację kwantyfikatorów, niezależnie
od tego, czy ustalono charakterystykę.
\end{twierdzenie}
\begin{proof}
Na mocy lematu \ref{lem:mod-qe-exists} wystarczy
wyeliminować $Y$ z formuły
\[
\exists Y\,
\left(\bigwedge_{i=1}^{r}P_i(Y,\bar x)=0
\land \bigwedge_{j=1}^{s}Q_j(Y,\bar x)\ne0\right),
\]
gdzie wszystkie $P_i,Q_j$ są wielomianami o współczynnikach
całkowitych. Niech $Q=\prod_{j=1}^{s}Q_j$ (dla $s=0$
przyjmujemy $Q=1$). W ciele warunek wszystkich
nierówności jest równoważny $Q\ne0$, ponieważ nie ma
dzielników zera.

Ustalmy ciało algebraicznie domknięte $K$ i wartości
parametrów $\bar a$. Ciało $K$ jest nieskończone:
gdyby było skończone, wielomian
$\prod_{c\in K}(Y-c)-1$ byłby niestały i nie miałby
pierwiastka w $K$, sprzecznie z algebraicznym domknięciem.
Jeśli wszystkie $P_i(Y,\bar a)$ są wielomianami zerowymi
(w tym gdy $r=0$), świadek istnieje dokładnie wtedy,
gdy $Q(Y,\bar a)$ nie jest wielomianem zerowym.
Rzeczywiście niezerowy wielomian ma co najwyżej tyle
pierwiastków, ile wynosi jego stopień: po znalezieniu
pierwiastka $c$ dzielimy przez $Y-c$ i stosujemy
indukcję po stopniu. W nieskończonym ciele można
zatem wybrać punkt spoza skończonego zbioru pierwiastków.

W pozostałym przypadku oznaczmy przez $G$ największy
wspólny dzielnik niezerowych $P_i(Y,\bar a)$.
Wspólne pierwiastki wielomianów $P_i$ to dokładnie
pierwiastki $G$: jeśli $c$ jest wspólnym pierwiastkiem,
$Y-c$ dzieli każdy $P_i$, więc dzieli $G$; odwrotnie
$G$ dzieli każdy $P_i$.
Jeżeli $G$ jest stały, świadek nie istnieje.
Jeżeli $\deg G>0$, niech $D\geq1$ będzie ustalonym
ograniczeniem stopni wszystkich wyjściowych $P_i$
w zmiennej $Y$. Wtedy $\deg G\leq D$.
Twierdzimy, że
\[
\exists c\in K\ (G(c)=0\land Q(c,\bar a)\ne0)
\quad\Longleftrightarrow\quad
G\nmid Q(Y,\bar a)^D .
\]
Jeżeli istnieje taki $c$, podzielność byłaby sprzeczna
z podstawieniem $c$. Jeżeli nie istnieje, każdy
pierwiastek $G$ jest pierwiastkiem $Q$.
Gdy $Q=0$, podzielność zachodzi automatycznie.
W przeciwnym razie rozkładamy $G$ na czynniki liniowe
w $K[Y]$, co wynika przez indukcję po stopniu
z algebraicznego domknięcia i dzielenia przez $Y-c$.
Każdy czynnik ma krotność co najwyżej $D$ w $G$,
a w $Q^D$ ma krotność co najmniej $D$.
Stąd $G$ dzieli $Q^D$.

Otrzymaliśmy skończoną alternatywę warunków:
wszystkie $P_i$ zerowe i $Q$ niezerowy; albo nie
wszystkie $P_i$ zerowe, $\deg G>0$ i
$G\nmid Q^D$. Są to warunki bez kwantyfikatorów
na współczynniki według lematu
\ref{lem:mod-qe-polynomial-branches}.
Równoważnik jest jednolity dla wszystkich ciał
algebraicznie domkniętych i nie wprowadza nowych
parametrów. Lemat \ref{lem:mod-qe-exists} kończy dowód.
\end{proof}

\begin{wniosek}[Zupełność w ustalonej charakterystyce]
\label{cor:mod-qe-acf-complete}
Teoria ciał algebraicznie domkniętych ustalonej
charakterystyki, dla której istnieje model, jest zupełna.
\end{wniosek}
\begin{proof}
Term domknięty w języku pierścieni reprezentuje
liczbę całkowitą interpretowaną w ciele przez wielokrotność
jedynki i ewentualny minus. Atom domknięty jest więc
równością takich liczb, a jego prawdziwość zależy
tylko od charakterystyki: w charakterystyce zero
różnica musi być zerowa w $\mathbb Z$, w charakterystyce
$p>0$ musi być podzielna przez $p$. Tak samo kombinacje
boolowskie atomów domkniętych zależą jedynie od
charakterystyki. Z twierdzenia \ref{thm:mod-qe-acf}
i wniosku \ref{cor:mod-qe-consequences} wynika zupełność.
\end{proof}

\begin{przyklad}[Eliminacja w konkretnej formule]
\label{ex:mod-qe-acf-example}
W każdym ciele algebraicznie domkniętym formuła
$\exists y\,(ay=1\land y\ne b)$ jest równoważna
$a\ne0\land ab\ne1$.
Jeśli $a=0$, pierwsze równanie nie ma rozwiązania.
Jeśli $a\ne0$, jedynym rozwiązaniem jest $y=a^{-1}$;
nierówność $y\ne b$ jest równoważna $ab\ne1$,
ponieważ mnożenie przez $a$ jest różnowartościowe.
W tym przykładzie wystarcza nawet dowolne ciało.
\end{przyklad}

\section{Ciała rzeczywiście domknięte}
Przechodzimy do języka uporządkowanych pierścieni
$\{0,1,+,-,\cdot,<\}$. W przeciwieństwie do przypadku ciał
algebraicznie domkniętych trzeba uwzględnić nie tylko
pierwiastki wielomianów, lecz także ich znaki pomiędzy
pierwiastkami. Przedstawimy algebraiczny dowód eliminacji,
działający również w ciałach niearchimedesowych.

\begin{definicja}[Ciało rzeczywiście domknięte]
\label{def:mod-qe-rcf}
Uporządkowane ciało $K$ nazywamy rzeczywiście domkniętym,
jeżeli $K(i)$, gdzie $i^2=-1$, jest algebraicznie domknięte.
Przez $T_{\mathrm{RCF}}$ oznaczamy teorię tych ciał
w języku uporządkowanych pierścieni, to znaczy zbiór
zdań prawdziwych we wszystkich takich ciałach.
\end{definicja}

Przyjmujemy powyższą charakterystykę jako definicję.
Ponieważ w ciele uporządkowanym kwadraty są nieujemne,
$Y^2+1$ nie ma pierwiastka w $K$ i jest nierozkładalny
jako wielomian stopnia drugiego. Dlatego $K(i)$ jest
ciałem, a każdy jego element ma jednoznaczną postać
$a+bi$ dla $a,b\in K$. Sprzężenie $a+bi\mapsto a-bi$
jest automorfizmem ustalającym $K$, co sprawdzamy
bezpośrednio z reguł dodawania i mnożenia oraz $i^2=-1$.
Przykładem ciała rzeczywiście domkniętego jest $\R$:
$\R(i)=\mathbb C$, a algebraiczne domknięcie $\mathbb C$
wynika z zasadniczego twierdzenia algebry.

\begin{lemat}[Znaki wielomianów]
\label{lem:mod-qe-rcf-signs}
Każdy niezerowy wielomian nad ciałem rzeczywiście
domkniętym rozkłada się na stałą niezerową oraz czynniki
liniowe i czynniki postaci $(Y-a)^2+b^2$, gdzie
$a,b\in K$, $b\ne0$. W szczególności jego znak
jest stały na każdym przedziale pomiędzy kolejnymi
różnymi pierwiastkami w $K$ oraz na przedziałach
nieograniczonych przed pierwszym i za ostatnim
pierwiastkiem. Jeśli nie ma pierwiastków w $K$,
jego znak jest stały na całym $K$.
\end{lemat}
\begin{proof}
Wielomian rozkłada się w $K(i)[Y]$ na czynniki liniowe
przez indukcję po stopniu, ponieważ $K(i)$ jest algebraicznie
domknięte. Sprzężenie zachowuje współczynniki wielomianu,
więc pierwiastki spoza $K$ występują parami sprzężonymi
z tymi samymi krotnościami. Krotność zachowuje się,
gdyż sprzężenie przenosi podzielność przez
$(Y-a-bi)^m$ na podzielność przez $(Y-a+bi)^m$.
Iloczyn pary czynników to $(Y-a)^2+b^2>0$ dla
każdego $Y\in K$. Pozostałe czynniki liniowe mają znak
ustalony przez położenie argumentu względem ich
pierwiastków. Iloczyn ma więc stały znak na każdym
z wymienionych przedziałów.
\end{proof}

Piszemy $\operatorname{sgn}(u)\in\{-1,0,1\}$ dla znaku
elementu $u$ ciała uporządkowanego. Znak niezerowego
wielomianu $H$ dla dostatecznie dużych argumentów
dodatnich jest znakiem jego współczynnika wiodącego,
a dla dostatecznie małych argumentów ujemnych jest
tym znakiem pomnożonym przez $(-1)^{\deg H}$.
W ciele rzeczywiście domkniętym wynika to z poprzedniego
lematu: po przekroczeniu wszystkich pierwiastków
czynniki liniowe są dodatnie, przed wszystkimi są
ujemne, a czynniki stopnia drugiego są dodatnie.
Liczba czynników liniowych z krotnościami ma tę samą
parzystość co stopień wielomianu.
Oznaczamy te dwa znaki przez $\operatorname{sgn}H(+\infty)$
i $\operatorname{sgn}H(-\infty)$; nie oznacza to
podstawiania nieskończoności za element ciała.

\begin{lemat}[Suma znaków na pierwiastkach]
\label{lem:mod-qe-sturm-tarski}
Niech $A\in K[Y]$ będzie niezerowym wielomianem
bez wielokrotnych pierwiastków (w $K(i)$), a $H\in K[Y]$.
Liczbę całkowitą
\[
\operatorname{Ta}(H,A)=
\sum_{\substack{r\in K\\ A(r)=0}}\operatorname{sgn}H(r)
\]
można wyznaczyć za pomocą działań na współczynnikach,
testów ich znaków i algorytmu Euklidesa, bez wyznaczania
samych pierwiastków.
\end{lemat}
\begin{proof}
Jeżeli $A$ jest stały, suma jest zerowa. W przeciwnym
razie niech $B$ będzie resztą z dzielenia $A'H$ przez $A$,
gdzie $A'$ oznacza formalną pochodną. Jeżeli $B=0$,
to $A'(r)H(r)=0$ dla każdego pierwiastka $r\in K$.
Ponieważ $A$ nie ma pierwiastków wielokrotnych, $A'(r)\ne0$:
pisząc $A=(Y-r)A_1$ otrzymujemy $A'(r)=A_1(r)\ne0$.
Stąd suma jest zerowa również w tym przypadku.

Dla $B\ne0$ wybieramy największy wspólny dzielnik
$G$ wielomianów $A,B$ i kładziemy $C=A/G$, $E=B/G$.
Są to wielomiany względnie pierwsze, a $\deg E<\deg C$.
W szczególności $C$ jest niestały. Pierwiastki $A$
usunięte przez dzielenie przez $G$ to dokładnie te
pierwiastki $A$ w $K$, dla których $H(r)=0$:
wspólny pierwiastek $A,B$ jest pierwiastkiem $G$
na mocy własności największego wspólnego dzielnika,
a przy pierwiastku $A$ zachodzi $B(r)=A'(r)H(r)$.
Dla pozostałych pierwiastków $r$ mamy $G(r)\ne0$,
$A'(r)=G(r)C'(r)$ i $B(r)=G(r)E(r)$. Zatem
\[
\frac{E(r)}{C'(r)}=H(r).
\]
Obie strony są niezerowe, ponieważ $C,E$ są względnie
pierwsze, a $C$ nie ma wielokrotnych pierwiastków.

Tworzymy ciąg wielomianów
$C_0=C$, $C_1=E$, a następnie
\[
C_{j+1}=-\operatorname{rem}(C_{j-1},C_j)
\]
aż do ostatniego niezerowego wielomianu $C_s$.
Tutaj $\operatorname{rem}$ oznacza resztę z dzielenia
w $K[Y]$. Stopnie maleją, więc procedura się kończy;
ostatni wielomian jest stały i niezerowy, ponieważ
algorytm Euklidesa zachowuje wspólne dzielniki
(por. lemat \ref{lem:mod-qe-polynomial-branches})
i $C,E$ są względnie pierwsze.

Dla skończonego ciągu znaków definiujemy liczbę zmian
znaku $V$ jako liczbę sąsiednich par różnych znaków
po usunięciu zer. Niech $V(-\infty)$ i $V(+\infty)$
oznaczają tę liczbę dla ciągu znaków wielomianów
$C_0,\ldots,C_s$ odpowiednio po lewej i po prawej
stronie wszystkich ich pierwiastków w $K$.
Wykażemy
\[
\operatorname{Ta}(H,A)=V(-\infty)-V(+\infty).
\]

Na każdym przedziale bez pierwiastków któregokolwiek
$C_j$ liczba zmian znaku jest stała na mocy lematu
\ref{lem:mod-qe-rcf-signs}.
Sąsiednie wielomiany $C_{j-1},C_j$ nie mają wspólnego
pierwiastka: wspólny pierwiastek przenosiłby się
przez równania dzielenia z resztą aż do ostatniej
stałej niezerowej. Jeżeli $C_j(r)=0$ dla $0<j<s$,
to z równania
$C_{j-1}=Q_jC_j-C_{j+1}$ wynika
$C_{j-1}(r)=-C_{j+1}(r)\ne0$.
Bezpośrednio po obu stronach $r$ znaki sąsiadów
są więc przeciwne i nie zmieniają się. W dwóch
parach sąsiednich z udziałem $C_j$ występuje
dokładnie jedna zmiana znaku niezależnie od znaku
$C_j$. Taki pierwiastek nie wnosi wkładu do różnicy
liczb zmian znaku. Jeżeli kilka wielomianów zeruje
się w $r$, ich indeksy nie są sąsiednie, więc
rozpatrywane pary są różne i można sumować wkłady.

Pozostaje pierwiastek $r$ wielomianu $C_0=C$.
Jest prosty. Z rozkładu $C=(Y-r)C_*$ i
$C_*(r)=C'(r)\ne0$ wynika, że znak $C$ tuż na lewo
od $r$ to $-\operatorname{sgn}C'(r)$, a tuż na prawo
to $\operatorname{sgn}C'(r)$. Znak $E$ pozostaje
$\operatorname{sgn}E(r)$. Zatem liczba zmian znaku
w pierwszej parze po lewej minus liczba zmian
po prawej wynosi $1$, jeżeli $C'(r)$ i $E(r)$ mają
ten sam znak, oraz $-1$ w przeciwnym razie.
Jest to $\operatorname{sgn}(E(r)/C'(r))
=\operatorname{sgn}H(r)$.
Pozostałe pary dają wkład zerowy, jak wyżej.
Sumując różnice po wszystkich pierwiastkach ciągu
w ich rosnącym porządku, otrzymujemy żądaną równość.

Wartości $V(\pm\infty)$ zależą wyłącznie od stopni
i znaków współczynników wiodących, według reguły
podanej przed lematem. Wszystkie użyte operacje
są zatem operacjami na współczynnikach.
\end{proof}

\begin{lemat}[Wyznaczanie układów znaków na pierwiastkach]
\label{lem:mod-qe-rcf-sign-determination}
Niech $A$ będzie jak w lemacie \ref{lem:mod-qe-sturm-tarski},
a $H_1,\ldots,H_m$ dowolnymi wielomianami w $K[Y]$.
Dla każdego $\sigma\in\{-1,0,1\}^m$ można obliczyć
liczbę pierwiastków $r\in K$ wielomianu $A$ takich, że
$\operatorname{sgn}H_j(r)=\sigma_j$ dla $1\leq j\leq m$,
korzystając wyłącznie z operacji i testów opisanych
w poprzednim lemacie.
\end{lemat}
\begin{proof}
Dla $v\in\{-1,0,1\}$ funkcje
\[
I_0(v)=1-v^2,\qquad
I_1(v)=\frac{v^2+v}{2},\qquad
I_{-1}(v)=\frac{v^2-v}{2}
\]
mają wartość $1$ dokładnie wtedy, gdy $v$ jest
odpowiednio równy $0,1,-1$, a w pozostałych
przypadkach mają wartość $0$, co sprawdzamy
przez podstawienie trzech możliwych wartości.
Żądana liczba pierwiastków jest więc równa
\[
\sum_{A(r)=0,\ r\in K}
\prod_{j=1}^m I_{\sigma_j}(\operatorname{sgn}H_j(r)).
\]
Po rozwinięciu iloczynu jest to skończona kombinacja
liniowa o współczynnikach wymiernych sum
\[
\operatorname{Ta}\left(\prod_{j=1}^m H_j^{e_j},A\right),
\qquad e_j\in\{0,1,2\}.
\]
Wynika to z multiplikatywności znaku, przy konwencji
$H_j^0=1$, także gdy $H_j(r)=0$.
Każdą sumę obliczamy poprzednim lematem.
Arytmetykę tych całkowitych wyników wykonujemy poza
rozważanym ciałem, w $\Q$; otrzymana liczba jest
nieujemną liczbą całkowitą, ponieważ liczy pierwiastki.
\end{proof}

\begin{twierdzenie}[Eliminacja kwantyfikatorów Tarskiego]
\label{thm:mod-qe-rcf}
Teoria ciał rzeczywiście domkniętych ma eliminację
kwantyfikatorów w języku uporządkowanych pierścieni.
Dokładniej, każda formuła tego języka ma równoważnik
bez kwantyfikatorów jednolity dla wszystkich ciał
rzeczywiście domkniętych.
\end{twierdzenie}
\begin{proof}
Na mocy lematu \ref{lem:mod-qe-exists} wystarczy
wyeliminować $Y$ z formuły $\exists Y\,\theta(Y,\bar x)$,
gdzie $\theta$ nie zawiera kwantyfikatorów.
Atomy są równościami lub nierównościami termów,
a termy reprezentują wielomiany o współczynnikach
całkowitych, jak w przypadku ciał algebraicznie domkniętych.
Po przeniesieniu termów na jedną stronę prawdziwość
$\theta$ zależy wyłącznie od znaków pewnej skończonej
listy wielomianów $P_1(Y,\bar x),\ldots,P_t(Y,\bar x)$.
Wystarczy rozstrzygnąć, które układy znaków tej listy
występują dla jakiegoś $Y\in K$, przy ustalonych
wartościach parametrów w ciele rzeczywiście domkniętym $K$.

Rozróżniamy wielomiany zerowe oraz rzeczywiste stopnie
pozostałych. Wielomiany zerowe mają zawsze znak $0$.
Niech $P$ będzie iloczynem pozostałych wielomianów
(pusty iloczyn wynosi $1$). Jeśli $P$ jest stały,
wszystkie znaki są stałe i wystarczy sprawdzić znaki
współczynników stałych. W przeciwnym razie bierzemy
\[
A=P/\gcd(P,P').
\]
Ponieważ każde ciało uporządkowane ma charakterystykę
zero (sumy dodatnich jedynek są dodatnie), $A$ ma
dokładnie te same pierwiastki co $P$ w $K(i)$,
ale wszystkie są proste. Rzeczywiście jeśli
$P=(Y-r)^d U$ w $K(i)[Y]$, $U(r)\ne0$, to
\[
P'=(Y-r)^{d-1}\bigl(dU+(Y-r)U'\bigr),
\]
a czynnik w nawiasie nie zeruje się w $r$.
Największy wspólny dzielnik ma więc w $r$ krotność
$d-1$. Algorytm Euklidesa nad $K$ daje ten sam
dzielnik (z dokładnością do niezerowej stałej) po
rozszerzeniu współczynników do $K(i)$, ponieważ
kolejne dzielenia i reszty nie ulegają zmianie.

Wszystkie układy znaków występujące w $K$ można
odczytać z następujących danych:
\begin{enumerate}
\item układu znaków na lewo od wszystkich pierwiastków $A$;
\item układów znaków w samych pierwiastkach $A$ w $K$;
\item układów znaków bezpośrednio na prawo od każdego
pierwiastka $A$ w $K$.
\end{enumerate}
Istotnie zbiór pierwiastków $A$ w $K$ jest skończony
i porządkuje się liniowo. Dopełnienie składa się
z przedziałów, na których wszystkie znaki są stałe
(lemat \ref{lem:mod-qe-rcf-signs}).
Każdy z tych przedziałów jest niepusty: między
dwoma różnymi elementami leży ich średnia, a z obu
stron skończonej listy można wybrać element przez
odjęcie lub dodanie $1$. Jeżeli nie ma pierwiastków
w $K$, wystarczy pierwszy układ znaków.

Pierwszy układ obliczamy ze stopni i współczynników
wiodących. Dla pozostałych dwóch stosujemy lemat
\ref{lem:mod-qe-rcf-sign-determination} do listy
wszystkich formalnych pochodnych $P_j^{(k)}$,
$0\leq k\leq\deg P_j$, dla niezerowych $P_j$.
Pochodna rzędu $0$ to sam wielomian.
Każdy układ znaków tych pochodnych, dla którego
obliczona liczba pierwiastków jest dodatnia, daje
układ znaków $P_j$ w pewnym pierwiastku $r$.
Daje też układ tuż na prawo od $r$: dla każdego
$P_j$ bierzemy znak pierwszej niezerowej pochodnej
$P_j^{(k)}(r)$, poczynając od $k=0$.
Istotnie rozwinięcie wielomianu w potęgach $Y-r$
ma współczynniki $P_j^{(k)}(r)/k!$, co wynika
przez różniczkowanie potęg $Y-r$ i podstawienie $Y=r$.
Jeżeli $k$ jest pierwszym niezerowym indeksem, to
$P_j=(Y-r)^k U$ z $U(r)=P_j^{(k)}(r)/k!\ne0$.
Dla $Y>r$ dostatecznie bliskiego $r$ znak $U(Y)$
jest znakiem $U(r)$, na mocy rozkładu z lematu
\ref{lem:mod-qe-rcf-signs}: wystarczy nie przekroczyć
żadnego pierwiastka $U$. Ponieważ $k!>0$, jest to
znak wskazanej pochodnej. Wspólny przedział po prawej
można wybrać dla całej skończonej listy; jego
prawy koniec można przyjąć jako następny pierwiastek
$P$ lub dowolny większy element, gdy następnego
pierwiastka nie ma. Wielomian zerowy zachowuje znak $0$.

Wyjaśnijmy, dlaczego otrzymana procedura daje formułę
bez kwantyfikatorów z parametrami $\bar x$, a nie tylko
test dla ustalonych parametrów. Początkowe stopnie
są ograniczone przez stopnie wyjściowych wielomianów.
Iloczyn, pochodne, największy wspólny dzielnik,
ilorazy i ciągi reszt mają stopnie ograniczone
przez liczby zależne wyłącznie od wyjściowej listy
i skończonej liczby wykonywanych operacji.
Algorytmy dzielenia i obliczania reszt mają
skończone drzewa rozgałęzień według rzeczywistych
stopni, jak w lemacie \ref{lem:mod-qe-polynomial-branches}.
Na każdej gałęzi wszystkie współczynniki są ilorazami
wielomianów w $\bar x$ o niezerowych mianownikach.
Znak ilorazu $U/V$ jest znakiem iloczynu $UV$,
gdyż $V^2>0$. Zatem wszystkie testy zerowości i znaku
są warunkami wielomianowymi na parametry. Każdy
test ma skończenie wiele wyników, a stopnie i liczba
operacji ograniczają długość gałęzi. Liczba wywołań
lematu \ref{lem:mod-qe-sturm-tarski} jest skończona,
bo lista pochodnych i zbiór wektorów wykładników
$\{0,1,2\}^m$ są skończone. Na każdym końcowym
liściu wszystkie liczby zmian znaków i liczby
pierwiastków są już ustalonymi liczbami całkowitymi.

Wybieramy te liście skończonego drzewa, na których
co najmniej jeden z obliczonych układów znaków
spełnia kombinację boolowską opisującą $\theta$.
Każdy liść jest koniunkcją warunków bez kwantyfikatorów,
a alternatywa wybranych liści jest żądanym równoważnikiem
$\exists Y\,\theta(Y,\bar x)$. Nie użyto żadnej
własności archimedesowej ani zupełności porządku.
Lemat \ref{lem:mod-qe-exists} kończy dowód.
\end{proof}

\begin{wniosek}[Zupełność i rozstrzygalność]
\label{cor:mod-qe-rcf-complete}
Teoria $T_{\mathrm{RCF}}$ jest zupełna i rozstrzygalna.
\end{wniosek}
\begin{proof}
Każdy term domknięty reprezentuje liczbę całkowitą,
a jej znak w ciele uporządkowanym jest zgodny ze
znakiem w $\mathbb Z$. Wszystkie atomy domknięte
mają więc ustaloną wartość logiczną, tak samo ich
kombinacje boolowskie. Twierdzenie \ref{thm:mod-qe-rcf}
i wniosek \ref{cor:mod-qe-consequences} dają zupełność.
Dowód eliminacji jest efektywny: rozgałęzienia według
stopni, rachunek wielomianowy, rozwinięcia iloczynów
i liczenie zmian znaków są skończonymi procedurami
na danych syntaktycznych. Po eliminacji kwantyfikatorów
z danego zdania obliczamy wartości jego domkniętych
termów jako liczb całkowitych i sprawdzamy atomy,
a następnie spójniki. Wynik rozstrzyga przynależność
zdania do $T_{\mathrm{RCF}}$.
\end{proof}

\section{Zbiory semialgebraiczne i rzuty}
\begin{definicja}[Zbiór semialgebraiczny]
\label{def:mod-qe-semialgebraic}
Niech $K$ będzie ciałem rzeczywiście domkniętym.
Zbiór $S\subseteq K^n$ nazywamy semialgebraicznym,
jeżeli jest skończoną kombinacją boolowską zbiorów
postaci $\{\bar x:P(\bar x)=0\}$ oraz
$\{\bar x:P(\bar x)>0\}$, gdzie $P\in K[X_1,\ldots,X_n]$.
\end{definicja}

Zbiory semialgebraiczne są dokładnie zbiorami
definiowalnymi z parametrami przez formuły bez
kwantyfikatorów w języku uporządkowanych pierścieni:
współczynniki wielomianów traktujemy jako parametry,
a w przeciwną stronę zapisujemy każdy atom jako
warunek na wielomian. Z twierdzenia \ref{thm:mod-qe-rcf}
i wniosku \ref{cor:mod-qe-definability} wynika zatem,
że każdy zbiór definiowalny z parametrami w tym języku
jest semialgebraiczny.

\begin{twierdzenie}[Twierdzenie Tarskiego--Seidenberga]
\label{thm:mod-qe-tarski-seidenberg}
Jeżeli $S\subseteq K^{n+m}$ jest semialgebraiczny,
to jego rzut na pierwsze $n$ współrzędnych jest
semialgebraiczny.
\end{twierdzenie}
\begin{proof}
Niech $\theta(\bar x,\bar y,\bar a)$ będzie formułą
bez kwantyfikatorów definiującą $S$, gdzie $\bar a$
jest skończoną krotką współczynników wielomianów.
Rzut jest zbiorem punktów $\bar x$, dla których
zachodzi $\exists\bar y\,\theta(\bar x,\bar y,\bar a)$.
Traktujemy $\bar a$ jako podstawienie dodatkowych
zmiennych wolnych i stosujemy twierdzenie
\ref{thm:mod-qe-rcf}. Otrzymujemy formułę bez
kwantyfikatorów, po podstawieniu parametrów
definiującą dokładnie ten rzut. Z definicji
\ref{def:mod-qe-semialgebraic} jest on semialgebraiczny.
\end{proof}

\begin{przyklad}[Rzut wykresu paraboli]
\label{ex:mod-qe-rcf-square}
W każdym ciele rzeczywiście domkniętym
$\exists y\,(y^2=x)$ jest równoważne $x\geq0$.
Konieczność wynika z uporządkowania ciała.
Dla $x=0$ świadkiem jest $0$; gdy $x>0$, wielomian
$Y^2-x$ ma pierwiastek $a+bi$ w $K(i)$.
Z równości $(a+bi)^2=x$ wynika
$2ab=0$ oraz $a^2-b^2=x$. Gdyby $a=0$, mielibyśmy
$x=-b^2\leq0$, sprzecznie z założeniem. Zatem $b=0$
i $a^2=x$. Jest to konkretny przypadek rzutu
semialgebraicznego.
\end{przyklad}

\begin{figure}[htbp]
\centering
\begin{tikzpicture}[x=1.25cm,y=1cm,>=Stealth]
\draw[->] (-.7,0)--(4.5,0) node[right] {$x$};
\draw[->] (0,-2.25)--(0,2.4) node[above] {$y$};
\draw[blue!65!black,thick,domain=-2:2,samples=81,variable=\t]
 plot ({\t*\t},{\t});
\draw[orange!85!black,very thick,->] (0,-2.85)--(4.4,-2.85);
\fill[orange!85!black] (0,-2.85) circle (2.5pt);
\node[below] at (2,-2.9) {$\{x:x\geq0\}$};
\foreach \x/\y in {1/1,2.25/1.5,4/2}{
 \draw[densely dashed,gray,->] (\x,\y)--(\x,-2.7);
 \fill[blue!65!black] (\x,\y) circle (1.7pt);
}
\node[blue!65!black] at (2.55,-1) {$y^2=x$};
\end{tikzpicture}
\caption{Rzut paraboli $y^2=x$ na oś $x$ w $\R^2$.
Kwantyfikator egzystencjalny odpowiada rzutowi.}
\label{fig:mod-qe-rcf-projection}
\end{figure}

Rysunek \ref{fig:mod-qe-rcf-projection} przedstawia
geometryczny sens przykładu \ref{ex:mod-qe-rcf-square}.

\begin{wniosek}[Zbiory definiowalne w jednej zmiennej]
\label{cor:mod-qe-rcf-one-variable}
Każdy zbiór definiowalny w $K$ z parametrami w języku
uporządkowanych pierścieni jest skończoną sumą punktów
i otwartych przedziałów, w tym przedziałów nieograniczonych.
\end{wniosek}
\begin{proof}
Eliminacja kwantyfikatorów daje kombinację warunków
znakowych dla skończonej listy wielomianów w jednej
zmiennej. Odrzucamy wielomiany zerowe, których znak
jest ustalony, i porządkujemy wszystkie różne pierwiastki
pozostałych wielomianów w $K$. Każdy z rozważanych
warunków ma stałą wartość logiczną na każdym otwartym
przedziale wyznaczonym przez tę listę (lemat
\ref{lem:mod-qe-rcf-signs}), a punkty listy rozpatrujemy
osobno. Zbiór jest sumą tych przedziałów i punktów,
na których kombinacja jest prawdziwa. W przypadku
pustej listy pierwiastków zbiór jest pusty albo równy $K$.
\end{proof}

\section{Arytmetyka Presburgera}
W odróżnieniu od teorii pełnej arytmetyki rozważamy
wyłącznie dodawanie, porządek i stałe. Mnożenie przez
ustaloną liczbę całkowitą jest skrótem termu, nie zaś
dodatkowym dwuargumentowym działaniem.

\begin{definicja}[Język arytmetyki Presburgera]
\label{def:mod-qe-presburger}
Niech $\mathcal L_{\mathrm P}$ będzie językiem
$\{0,1,+,-,<,(\equiv_m)_{m\geq2}\}$, gdzie minus jest
jednoargumentowy, a $\equiv_m$ jest symbolem relacji
dwuargumentowej. W strukturze $\mathfrak Z$ o uniwersum
$\mathbb Z$ symbole arytmetyczne interpretujemy standardowo,
a relacje przez
\[
a\equiv_m b \quad\Longleftrightarrow\quad m\mid(a-b).
\]
Przez $T_{\mathrm P}$ oznaczamy teorię
$\operatorname{Th}(\mathfrak Z)$ w tym języku.
\end{definicja}

Dla $m\geq2$ kongruencję można zdefiniować w języku
bez symboli $\equiv_m$ formułą
$\exists z\,(a-b=mz)$. Dodanie tych symboli nie zwiększa
więc zbiorów definiowalnych w strukturze standardowej,
ale pozwala wyrazić kongruencje bez kwantyfikatorów.
Będziemy używać także zapisu $\equiv_1$, oznaczającego
zawsze prawdziwy warunek. Wszystkie liczby całkowite
w termach oznaczają wielokrotności $1$.

\begin{lemat}[Normalizacja ograniczeń jednej zmiennej]
\label{lem:mod-qe-presburger-normal}
Każdą formułę postaci $\exists y\,\theta(y,\bar x)$
z $\theta$ będącą skończoną koniunkcją literałów
w języku $\mathcal L_{\mathrm P}$ można przekształcić
w równoważną modulo $T_{\mathrm P}$ skończoną alternatywę
formuł postaci
\[
\begin{aligned}
\chi(\bar x)\land \exists z\,\bigg(&
\bigwedge_{i\in I}\ell_i(\bar x)\leq z
\land \bigwedge_{j\in J}z\leq u_j(\bar x)\\
&{}\land\bigwedge_{k\in K}
 z\equiv_{m_k}t_k(\bar x)\bigg),
\end{aligned}
\]
gdzie $\chi$ nie zawiera kwantyfikatorów, wszystkie
zbiory indeksów są skończone, a $\ell_i,u_j,t_k$
są termami języka $\mathcal L_{\mathrm P}$.
W zapisie $\leq$ używamy skrótu $\neg(>)$.
\end{lemat}
\begin{proof}
Wystarczy wykazać równoważność w $\mathfrak Z$ dla
każdego wartościowania parametrów, ponieważ jej
uniwersalne domknięcie należy wtedy do $T_{\mathrm P}$.

Każdy term jest wyrażeniem liniowym o współczynnikach
całkowitych: zmienna i stałe mają tę postać, a dodawanie
i negowanie zachowują ją. Przenosząc wyrażenia na
odpowiednie strony, sprowadzamy atomy do warunków
\[
ay<t,\qquad ay=t,\qquad ay\equiv_m t,
\]
gdzie $a\in\mathbb Z$ jest ustalone syntaktycznie,
a $t$ nie zawiera $y$. Negację ścisłej nierówności
zapisujemy jako nierówność słabą; negację równości
jako alternatywę dwóch ścisłych nierówności.
Równość zastępujemy dwiema nierównościami słabymi.
Nierówności ścisłe przekształcamy w słabe,
korzystając z równoważności $s<t$ i $s\leq t-1$ na
$\mathbb Z$. Negację kongruencji zapisujemy jako
\[
ay\not\equiv_m t
\quad\Longleftrightarrow\quad
\bigvee_{r=1}^{m-1} ay\equiv_m t+r,
\]
gdyż każda liczba całkowita ma dokładnie jedną
resztę w zbiorze $\{0,\ldots,m-1\}$ przy dzieleniu
przez $m$. Rozdzielając alternatywy, otrzymujemy
skończenie wiele koniunkcji. Warunki ze współczynnikiem
$a=0$ przenosimy do $\chi(\bar x)$.

Ustalmy jedną z pozostałych koniunkcji. Niech $L$
będzie wspólną dodatnią wielokrotnością wszystkich
niezerowych wartości $|a|$ (jeśli nie ma takich,
bierzemy $L=1$). Wprowadzamy nową zmienną $z=Ly$
i warunek $z\equiv_L0$. Każde ograniczenie z niezerowym
$a$ mnożymy przez dodatnią liczbę całkowitą $L/|a|$.
W nierównościach zachowuje to porządek; współczynnik
przy $y$ staje się $L$ lub $-L$. W kongruencjach
dla dodatniej liczby całkowitej $c=L/|a|$ stosujemy
równoważność
\[
ay\equiv_m t\quad\Longleftrightarrow\quad
cay\equiv_{cm}ct.
\]
Wynika ona z równoważności $m\mid v$ i $cm\mid cv$
dla $c>0$: można skrócić obie strony równości
$cv=cmw$ przez $c$ w $\mathbb Z$.
Po podstawieniu $z=Ly$ otrzymujemy współczynnik
$1$ lub $-1$ przy $z$. W przypadku $-1$ zmieniamy
znaki obu stron, odwracając kierunek nierówności
i zachowując kongruencję.
Wreszcie warunek $z\equiv_L0$ gwarantuje istnienie
$y\in\mathbb Z$ z $z=Ly$; dzięki temu podstawienie
jest odwracalne na zbiorze rozważanych świadków.
Dostajemy żądaną postać.
\end{proof}

\begin{twierdzenie}[Eliminacja kwantyfikatorów w arytmetyce Presburgera]
\label{thm:mod-qe-presburger}
Teoria $T_{\mathrm P}$ ma eliminację kwantyfikatorów
w języku $\mathcal L_{\mathrm P}$.
\end{twierdzenie}
\begin{proof}
Na mocy lematów \ref{lem:mod-qe-exists} oraz
\ref{lem:mod-qe-presburger-normal} wystarczy usunąć
kwantyfikator egzystencjalny z koniunkcji ograniczeń
dolnych, górnych i kongruencji w znormalizowanej zmiennej
$z$. Ponownie dowodzimy równoważności w $\mathbb Z$;
jej uniwersalne domknięcie należy do $T_{\mathrm P}$.

Niech $D$ będzie dodatnią wspólną wielokrotnością
wszystkich modułów $m_k$; dla pustej listy bierzemy
$D=1$. Warunki kongruencji są niezmiennicze
przy zmianie $z$ o dowolną wielokrotność $D$.

Jeśli istnieje co najmniej jedno ograniczenie dolne,
to rozdzielamy skończenie wiele przypadków według tego,
które $\ell_i$ jest największe. W przypadku indeksu
$i$ dołączamy koniunkcję
$\bigwedge_{h\in I}\ell_h\leq\ell_i$.
Oznaczmy $\ell=\ell_i$.
Wystarczy sprawdzić $D$ wartości
\[
z=\ell+r,\qquad r=0,\ldots,D-1.
\]
Istotnie, dla dowolnego świadka $z\geq\ell$
zapisujemy $z-\ell=qD+r$ z $q\geq0$ i
$0\leq r<D$. Liczba $\ell+r$ spełnia te same
kongruencje, leży nie niżej niż $\ell$ i nie wyżej
niż $z$, więc zachowuje wszystkie ograniczenia górne.
Implikacja odwrotna jest natychmiastowa: jeden
z testowanych kandydatów spełniający ograniczenia
jest świadkiem. Otrzymujemy formułę bez kwantyfikatorów
\[
\bigvee_{i\in I}
\left[
\bigwedge_{h\in I}\ell_h\leq\ell_i
\land
\bigvee_{r=0}^{D-1}
\left(
\bigwedge_{j\in J}\ell_i+r\leq u_j
\land
\bigwedge_{k\in K}\ell_i+r\equiv_{m_k}t_k
\right)
\right].
\]
Jeżeli $I=\varnothing$, lecz $J\ne\varnothing$,
analogicznie rozdzielamy przypadki według najmniejszego
ograniczenia górnego $u_j$ i testujemy $u_j-r$ dla
$r=0,\ldots,D-1$. Każdy świadek można przesunąć
w górę o wielokrotność $D$ do tego zbioru kandydatów,
nie przekraczając $u_j$. Gdy $I=J=\varnothing$,
testujemy $z=r$ dla $0\leq r<D$, ponieważ każda
liczba całkowita jest kongruentna modulo $D$ z dokładnie
jedną z tych liczb. W każdym przypadku otrzymujemy
skończoną kombinację boolowską atomów z parametrami
$\bar x$, bez nowych zmiennych wolnych.
\end{proof}

\begin{przyklad}[Eliminacja z ograniczeniami i kongruencją]
\label{ex:mod-qe-presburger-interval}
Rozważmy formułę
$\exists y\,(a\leq y\land y\leq b\land y\equiv_3 1)$.
Wybieramy $D=3$, a jedynym ograniczeniem dolnym jest $a$.
Z dowodu twierdzenia \ref{thm:mod-qe-presburger}
otrzymujemy równoważnik
\[
\bigvee_{r=0}^{2}
(a+r\leq b\land a+r\equiv_3 1).
\]
Na przykład dla $a=5$ i $b=7$ świadkiem jest $7$.
Dla $a=5$ i $b=6$ nie ma świadka: jedyny spośród
$5,6,7$ kandydat kongruentny z $1$ modulo $3$ to $7$,
który przekracza $b$.
\end{przyklad}

\begin{wniosek}[Rozstrzygalność]
\label{cor:mod-qe-presburger-decidable}
Teoria $T_{\mathrm P}$ jest rozstrzygalna, tj. istnieje
algorytm rozstrzygający, czy dane zdanie jej języka
należy do $T_{\mathrm P}$.
\end{wniosek}
\begin{proof}
Konstrukcje w lemacie \ref{lem:mod-qe-exists}, lemacie
\ref{lem:mod-qe-presburger-normal} i twierdzeniu
\ref{thm:mod-qe-presburger} są efektywne: rozpisanie
skończonych alternatyw, obliczanie współczynników
całkowitych, wspólnych wielokrotności oraz podstawienia
termów są skończonymi procedurami. Eliminujemy kwantyfikatory
od najgłębiej zagnieżdżonych podformuł. Wynikiem dla
zdania jest zdanie bez kwantyfikatorów. Obliczamy wartości
jego termów domkniętych w $\mathbb Z$, sprawdzamy
równości, nierówności i kongruencje oraz obliczamy
wartość kombinacji boolowskiej. Zdanie należy do
$T_{\mathrm P}$ dokładnie wtedy, gdy wynik jest prawdziwy.
Nie twierdzimy, że procedura ma małą złożoność.
\end{proof}

\begin{przyklad}[Kongruencje są potrzebne w tym języku]
\label{ex:mod-qe-presburger-no-congruence}
W języku $\{0,1,+,-,<\}$ teoria standardowej struktury
na $\mathbb Z$ nie ma eliminacji kwantyfikatorów.
Formuła $\exists y\,(x=y+y)$ definiuje zbiór liczb
parzystych. Każdy atom bez kwantyfikatorów w jednej
zmiennej można sprowadzić do równości lub nierówności
liniowej z całkowitymi współczynnikami.
Jeżeli współczynnik przy $x$ jest zerowy, wartość
logiczna atomu jest stała. Jeżeli jest niezerowy,
równość może zachodzić dla co najwyżej jednej
wartości $x$, a prawdziwość nierówności stabilizuje
się dla dostatecznie dużych dodatnich $x$ (znak
niezerowego współczynnika wyznacza ostateczny znak
wyrażenia liniowego).
Skończona kombinacja boolowska takich atomów również
ma stałą wartość logiczną dla wszystkich dostatecznie
dużych dodatnich $x$. Zbiór liczb parzystych nie ma
tej własności, ponieważ dowolnie duże liczby parzyste
i nieparzyste występują na przemian. Nie istnieje więc
równoważnik bez kwantyfikatorów w języku bez kongruencji.
\end{przyklad}

\begin{uwaga}[Wersja na liczbach naturalnych]
\label{rem:mod-qe-presburger-naturals}
Arytmetykę Presburgera często formułuje się na $\N$
w języku z dodawaniem, porządkiem, stałymi $0,1$
i kongruencjami, lecz bez jednoargumentowego minusa.
Zdania tej struktury tłumaczymy na zdania o $\mathbb Z$:
każde $\exists y\,\psi$ zastępujemy przez
$\exists y\,(0\leq y\land\psi)$, a każde
$\forall y\,\psi$ przez $\forall y\,(0\leq y\to\psi)$,
rekurencyjnie tłumacząc podformuły.
Indukcja po formułach pokazuje zgodność prawdziwości
przy nieujemnych wartościach zmiennych wolnych:
atomy i spójniki są zgodne, a kwantyfikatory
ograniczono dokładnie do $\N$ (przy konwencji $0\in\N$).
Eliminacja w $\mathbb Z$ daje formułę bez kwantyfikatorów,
której atomy można zapisać z termami bez minusa:
przenosimy wszystkie składniki o ujemnych
współczynnikach na drugą stronę równości, nierówności
lub kongruencji. Otrzymane współczynniki są nieujemne
i oznaczają iterowane dodawanie. Daje to eliminację
również w teorii standardowej struktury na $\N$
z wymienionymi symbolami.
\end{uwaga}

\section{Kiedy eliminacja nie zachodzi}
\begin{przyklad}[Porządki dyskretne i gęste]
\label{ex:mod-qe-not-linear-orders}
Teoria wszystkich porządków liniowych nie ma eliminacji
kwantyfikatorów w języku $\{<\}$. W strukturach
$(\mathbb Z,<)$ i $(\Q,<)$ krotki $(0,1)$ zgadzają się
co do równości i porządku, a zatem co do wszystkich
formuł bez kwantyfikatorów. Formuła
$\exists y\,(x<y\land y<z)$ jest fałszywa przy $(0,1)$
w pierwszej strukturze i prawdziwa w drugiej.
Kryterium \ref{thm:mod-qe-semantic} wyklucza eliminację.
\end{przyklad}

\begin{przyklad}[Zupełność nie wystarcza]
\label{ex:mod-qe-complete-not-qe}
Zupełna teoria $\operatorname{Th}(\mathbb Z,<)$
nie ma eliminacji kwantyfikatorów.
W jednym modelu $(\mathbb Z,<)$ krotki $(0,1)$
i $(0,2)$ mają ten sam typ bez kwantyfikatorów,
bo każda składa się z dwóch różnych rosnących elementów.
Formuła $\exists y\,(x<y\land y<z)$ rozróżnia te krotki:
pierwsza nie ma świadka, a druga ma świadka $1$.
Zatem warunek z twierdzenia \ref{thm:mod-qe-semantic}
nie zachodzi mimo zupełności teorii.
\end{przyklad}

Powyższe kryteria i przykłady pozwolą następnie badać
zupełność modelową, kategoryczność oraz własności
zbiorów definiowalnych.
